\documentclass[10pt]{article}
\usepackage[T1]{fontenc}
\usepackage[utf8]{inputenc}
\usepackage[french]{babel}
\usepackage{amsmath,amssymb,mathrsfs,mathpazo}
\usepackage{geometry}
\geometry{a4paper,margin=2cm}
\usepackage{graphicx,array,tabularx,tabularray,colortbl,multirow,multicol,booktabs,enumitem,pifont,ulem,fancybox,diagbox,cancel}
\usepackage{pstricks,pstricks-add,pst-plot,pst-tree,pst-node,pst-text,pst-func,pst-math,pst-bspline,pst-3dplot,pst-eucl,pst-solides3d}
\let\pscalculate\undefined
\usepackage{pst-bezier}
\usepackage{tikz,pgfplots,tkz-tab}
\usetikzlibrary{arrows,arrows.meta,patterns,calc,positioning,trees,decorations.markings,shapes.geometric,shapes.misc,shapes.arrows,intersections,decorations.pathmorphing,decorations.pathreplacing,backgrounds,shadows,babel}
\pgfplotsset{compat=1.18}
\usepgfplotslibrary{fillbetween}
\usepackage[np]{numprint}
\usepackage{listings,makecell,pas-tableur,twoopt}
\let\barre\undefined
\usepackage[pstricks]{tablvar}
\definecolor{darkgreen}{RGB}{0,100,0}
\providecommand{\red}{\color{red}}
\providecommand{\blue}{\color{blue}}
\providecommand{\green}{\color{green}}
\providecommand{\black}{\color{black}}
\providecommand{\white}{\color{white}}
\providecommand{\gray}{\color{gray}}
\providecommand{\cyan}{\color{cyan}}
\providecommand{\magenta}{\color{magenta}}
\providecommand{\yellow}{\color{yellow}}
\providecommand{\violet}{\color{violet}}
\providecommand{\orange}{\color{orange}}
\providecommand{\euro}{\texteuro}
\providecommand{\EUR}{\texteuro}
\providecommand{\ieme}{\textsuperscript{e}}
\providecommand{\degres}{\ensuremath{^\circ}}
\providecommand{\degree}{\ensuremath{^\circ}}
\providecommand{\vv}[1]{\overrightarrow{#1}}
\providecommand{\Coord}[1]{#1}
\providecommand{\decosix}{}
\providecommand{\starredbullet}{\ensuremath{\bullet}}
\pagestyle{empty}
\setlength{\parindent}{0pt}
\usepackage{adjustbox}\DeclareUnicodeCharacter{2236}{:}\DeclareUnicodeCharacter{2190}{\ensuremath{\leftarrow}}
\providecommand{\R}{\mathbb{R}}
\providecommand{\N}{\mathbb{N}}
\providecommand{\Z}{\mathbb{Z}}
\providecommand{\Q}{\mathbb{Q}}
\providecommand{\D}{\mathbb{D}}
\providecommand{\C}{\mathbb{C}}
\providecommand{\e}{{\mathrm{e}}}
\providecommand{\vect}[1]{\overrightarrow{#1}}
\providecommand{\vectt}[1]{\overrightarrow{\mathrm{#1}}}
\providecommand{\vvt}[1]{\overrightarrow{\mathrm{#1}}}
\providecommand{\np}[2][]{{#2~#1}}
\providecommand{\vv}[1]{{\overrightarrow{#1}}}
\renewcommand{\d}{{\mathrm{d}}}
\providecommand{\degres}{{^\circ}}
\providecommand{\degre}{{^\circ}}
\providecommand{\blue}{\color{blue}}
\providecommand{\red}{\color{red}}
\providecommand{\enskip}{\quad}
\providecommand{\llbracket}{\lbrack\!\lbrack}
\providecommand{\rrbracket}{\rbrack\!\rbrack}
\providecommand{\euro}{{\text{€}}}
\providecommand{\no}{{\text{n°}}}
\providecommand{\interval}[2]{{[#1;#2]}}
\providecommand{\intervFO}[2]{{[#1;#2[}}
\providecommand{\dotfill}{\ldots}

\providecommand{\ds}{\displaystyle}
\providecommand{\textding}[1]{#1}
\providecommand{\Oij}{\(\mathrm{O};\vec{\imath},\vec{\jmath}\)}
\providecommand{\Oijk}{\(\mathrm{O};\vec{\imath},\vec{\jmath},\vec{k}\)}
\providecommand{\Ouv}{\(\mathrm{O};\vec{u},\vec{v}\)}
\providecommand{\intervFF}[2]{\ensuremath{[#1;#2]}}
\renewcommand{\euro}{\ifmmode\text{\texteuro}\else\texteuro\fi}
\DeclareUnicodeCharacter{00B0}{\ifmmode^{\circ}\else\textdegree\fi}
\DeclareUnicodeCharacter{20AC}{\ifmmode\text{\texteuro}\else\texteuro\fi}
% Énoncés : Éducation nationale. Transcription et corrigés : APMEP, Denis Vergès et contributeurs. Découpage et classement : CoopMaths, Sébastien Lozano, Éric Elter et contributeurs. Adaptation de la présentation : Maths Explorer.
% CC BY-SA : https://coopmaths.fr/www/about/
% https://coopmaths.fr/alea/static/e3c/2020/tex/e3c_2020_00_sujet34_4_cor.tex
\begin{document}
\noindent\textit{APMEP / CoopMaths. CC BY-SA. Présentation adaptée par Maths Explorer.}\par\medskip
\(f(x) = x^3 - 6x^2 + 9x + 1\).

\subsection*{1.}

La fonction polynôme \(f\) est dérivable sur \(\mathbb{R}\) et sur \([-1\,;\,5]\), et sa dérivée est :
\[
f'(x) = 3x^2 - 12x + 9 = 3(x^2 - 4x + 3).
\]

\subsection*{2.}

Le nombre 1 est une racine évidente du trinôme \(x^2 - 4x + 3\), et comme le produit des racines est \(\dfrac{c}{a} = \dfrac{3}{1} = 3\), l'autre racine est 3.

On sait que \(x^2 - 4x + 3 = (x - 1)(x - 3)\), donc \(f'(x) = 3(x - 1)(x - 3)\).

\subsection*{3.}

Le trinôme est du signe de \(a = 1 > 0\), donc positif sauf sur l'intervalle \(\left]1\,;\,3\right[\).  

Il en résulte que sur l'intervalle \([-1\,;\,5]\) :
\[f(-1) = -1 - 6 - 9 + 1 = -15 \quad ; \quad f(1) = 1 - 6 + 9 + 1 = 5\]
\[f(3) = 27 - 54 + 27 + 1 = 1 \quad ; \quad f(5) = 125 - 150 + 45 + 1 = 21\]
\begin{center}
\begin{tikzpicture}
\tkzTabInit[lgt=2.5, espcl=2]{$x$ / 1, {Signe de $f'(x)$} / 1, {$f$} / 2}{${-1}$, ${1}$, ${3}$, ${5}$}
\tkzTabLine{,+,0,-,0,+,}
\tkzTabVar{-/{$-15$},+/{$5$},-/{$1$},+/{$21$}}{/}
\end{tikzpicture}
\end{center}

\subsection*{4.}

On sait qu'une équation de la tangente \(\mathcal{T}\) à la courbe représentative de la fonction \(f\) au point d'abscisse \(0\) est :
\[
M(x\,;\,y) \in \mathcal{T} \iff y - f(0) = f'(0)(x - 0).
\]
Avec \(f(0) = 1\) et \(f'(0) = 9\), on obtient :
\[
M(x\,;\,y) \in \mathcal{T} \iff y = 9x + 1.
\]

\subsection*{5.}

Si une autre tangente est parallèle à \(\mathcal{T}\), son coefficient directeur est égal à 9.

Or,
\begin{align*}
f'(x) &= 9 \\
3(x^2 - 4x + 3) &= 9 \\
x^2 - 4x + 3 &= 3 \\
x^2 - 4x &= 0 \\
x(x - 4) &= 0.
\end{align*}

On retrouve \(x = 0\) pour la tangente \(\mathcal{T}\), et \(x = 4\), pour lequel :
\[
f(4) = 64 - 96 + 36 + 1 = 5.
\]  
La tangente à la courbe représentative de la fonction \(f\) au point \(B(4\,;\,5)\) est parallèle à \(\mathcal{T}\).

\end{document}
